\chapter{ROS 2 の基本概念とコマンドラインツール}
\label{chap:ros2_concepts}

% この章で学ぶこと
\begin{goalbox}
  \begin{itemize}
    \item ROS 2 のプログラムの単位であるノード
    \item ノードどうしの通信の方法（トピック・サービス・アクション）と、その使い分け
    \item やり取りするデータの型（メッセージ型）の調べ方
    \item ノードの設定値であるパラメータ
    \item \cmd{ros2} コマンドと rqt を使って、動いている ROS 2 のシステムを調べる方法
  \end{itemize}
\end{goalbox}

\ref{chap:ros2_install}章では、turtlesim を動かして、ROS 2 のノードどうしが通信する様子を少しだけ見ました。
この章では、ROS 2 の基本となる考え方を、turtlesim を使いながら 1 つずつ学びます。
あわせて、動いている ROS 2 のシステムの中をのぞくための \cmd{ros2} コマンドの使い方を身につけます。
これらのコマンドは、自分でノードを書くようになってからも、動作の確認やトラブルの調査で毎日のように使います。

この章では、次の 2 つのノードを起動した状態で進めます。
それぞれ別のターミナル（またはタブ）で起動し、そのままにしておいてください。

\begin{terminal}[1 つ目のターミナル：turtlesim]
$ ros2 run turtlesim turtlesim_node
\end{terminal}

\begin{terminal}[2 つ目のターミナル：キーボード操作]
$ ros2 run turtlesim turtle_teleop_key
\end{terminal}

以降のコマンドは、3 つ目のターミナルで実行します。

\section{ノード}
\label{sec:ros2-concepts-node}

\subsection{ノードとは}

\textbf{ノード}は、ROS 2 のプログラムの単位です。
\ref{sec:what-is-ros-components}節で見たように、ロボットのソフトウェアは、「センサのデータを読む」「地図を作る」「モータを動かす」といった、1 つの役割を持つノードの集まりとして作ります。
基本的には、1 つのノードが 1 つのプロセス（\ref{sec:linux-internals-process}節）として動きます。

\subsection{ノードを調べる}

動いているノードの一覧は、\cmd{ros2 node list} で表示できます。

\begin{terminal}[ノードの一覧を表示する]
$ ros2 node list
/teleop_turtle
/turtlesim
\end{terminal}

ノードの名前は \cmd{/} から始まります。
\cmd{ros2 node info} を使うと、そのノードがどのトピックやサービスを使っているかがわかります。

\begin{terminal}[ノードの詳しい情報を表示する]
$ ros2 node info /turtlesim
/turtlesim
  Subscribers:
    /parameter_events: rcl_interfaces/msg/ParameterEvent
    /turtle1/cmd_vel: geometry_msgs/msg/Twist
  Publishers:
    /parameter_events: rcl_interfaces/msg/ParameterEvent
    /rosout: rcl_interfaces/msg/Log
    /turtle1/color_sensor: turtlesim/msg/Color
    /turtle1/pose: turtlesim/msg/Pose
  Service Servers:
    /clear: std_srvs/srv/Empty
    /kill: turtlesim/srv/Kill
    /reset: std_srvs/srv/Empty
    /spawn: turtlesim/srv/Spawn
    ...
  Action Servers:
    /turtle1/rotate_absolute: turtlesim/action/RotateAbsolute
  ...
\end{terminal}

\cmd{Subscribers} や \cmd{Publishers} などの意味は、この後の節で順に説明します。

\subsection{ノードの名前を変えて起動する}

同じノードを、別の名前で起動することもできます。
\cmd{ros2 run} の後ろに \cmd{--ros-args --remap \_\_node:=名前} を付けると、ノードの名前を変えられます。

\begin{terminal}[ノードの名前を変えて起動する]
$ ros2 run turtlesim turtlesim_node --ros-args --remap __node:=my_turtle
\end{terminal}

別のターミナルで \cmd{ros2 node list} を実行すると、\cmd{/my\_turtle} というノードが増えていることがわかります。
このように、起動するときの名前や設定を変える方法は、\ref{chap:launch}章で詳しく学びます。
確認したら、\key{Ctrl}+\key{C} で止めておきましょう。

\section{トピック}
\label{sec:ros2-concepts-topic}

\subsection{ノードどうしの 3 つの通信方法}

ROS 2 のノードどうしは、主に\textbf{トピック}・\textbf{サービス}・\textbf{アクション}の 3 つの方法で通信します（図\ref{fig:ros2-concepts-communication}）。
それぞれ得意な場面が違うので、目的に合わせて使い分けます。

\begin{figure}[tb]
  \centering
  % 図：トピック・サービス・アクションの通信の違い（chapters/ros2/ros2_concepts.tex）
\begin{tikzpicture}[
    node/.style={draw, ellipse, font=\footnotesize, fill=blue!8, minimum width=20mm, minimum height=8mm, inner sep=1pt},
    lab/.style={font=\small\bfseries, anchor=west},
    msg/.style={font=\scriptsize, midway, fill=white, inner sep=1pt}]
  % トピック
  \node[lab] at (-1.2,0.9) {トピック（出版・購読）};
  \node[node] (p) at (0,0) {パブリッシャ};
  \node[draw, fill=orange!10, font=\scriptsize\ttfamily, inner sep=3pt] (t) at (4.4,0) {/scan};
  \node[node] (s1) at (7.4,0.45) {サブスクライバ};
  \node[node] (s2) at (7.4,-0.45) {サブスクライバ};
  \draw[figarrow] (p) -- node[msg, above=1pt]{データを送り続ける} (t);
  \draw[figarrow] (t) -- (s1.west);
  \draw[figarrow] (t) -- (s2.west);
  % サービス
  \node[lab] at (-1.2,-1.5) {サービス（要求・応答）};
  \node[node] (c) at (0,-2.4) {クライアント};
  \node[node] (sv) at (7.4,-2.4) {サーバ};
  \draw[figarrow] ([yshift=1.5mm]c.east) -- node[msg, above=1pt]{要求（リクエスト）} ([yshift=1.5mm]sv.west);
  \draw[figarrow] ([yshift=-1.5mm]sv.west) -- node[msg, below=1pt]{応答（レスポンス）を 1 回返す} ([yshift=-1.5mm]c.east);
  % アクション
  \node[lab] at (-1.2,-3.9) {アクション（目標・途中経過・結果）};
  \node[node] (ac) at (0,-5.1) {クライアント};
  \node[node] (as) at (7.4,-5.1) {サーバ};
  \draw[figarrow] ([yshift=3mm]ac.east) -- node[msg, above=1pt]{目標（ゴール）} ([yshift=3mm]as.west);
  \draw[figarrow, dashed] (as.west) -- node[msg]{途中経過（フィードバック）を何度も} (ac.east);
  \draw[figarrow] ([yshift=-3mm]as.west) -- node[msg, below=1pt]{結果（リザルト）} ([yshift=-3mm]ac.east);
\end{tikzpicture}

  \caption{トピック・サービス・アクションの通信の違い}
  \label{fig:ros2-concepts-communication}
\end{figure}

\subsection{トピックとは}

\textbf{トピック}は、データを流すための名前の付いた通り道です。
データを送るノードを\textbf{パブリッシャ}（出版者）、受け取るノードを\textbf{サブスクライバ}（購読者）といいます。
パブリッシャがトピックにデータを\textbf{パブリッシュ}（出版）すると、そのトピックを\textbf{サブスクライブ}（購読）しているすべてのノードにデータが届きます。

雑誌の出版と購読にたとえるとわかりやすいでしょう。
出版社は、誰が読んでいるかを気にせずに雑誌を出版し続け、購読者は、購読している雑誌が届くのを待ちます。
同じように、パブリッシャとサブスクライバは、お互いのことを知らなくても、トピックの名前さえ同じなら通信できます。
1 つのトピックに、複数のパブリッシャやサブスクライバがいても構いません。

トピックは、センサのデータやロボットへの速度の指令のように、\textbf{次々と流れ続けるデータ}を送るのに向いています。
ROS 2 の通信の中で、最もよく使われる方法です。

\subsection{トピックを調べる}

トピックの一覧は、\cmd{ros2 topic list} で表示できます。
\cmd{-t} を付けると、トピックに流れるデータの型（メッセージ型）も表示されます。

\begin{terminal}[トピックの一覧を表示する]
$ ros2 topic list -t
/parameter_events [rcl_interfaces/msg/ParameterEvent]
/rosout [rcl_interfaces/msg/Log]
/turtle1/cmd_vel [geometry_msgs/msg/Twist]
/turtle1/color_sensor [turtlesim/msg/Color]
/turtle1/pose [turtlesim/msg/Pose]
\end{terminal}

\cmd{/turtle1/cmd\_vel} はカメへの速度の指令、\cmd{/turtle1/pose} はカメの今の位置と向きを流すトピックです。
\cmd{/parameter\_events} と \cmd{/rosout} は、ROS 2 が自動で用意するトピックです。

\cmd{ros2 topic echo} を使うと、トピックに流れているデータをそのまま表示できます。

\begin{terminal}[トピックに流れるデータを表示する]
$ ros2 topic echo /turtle1/pose
x: 5.544444561004639
y: 5.544444561004639
theta: 0.0
linear_velocity: 0.0
angular_velocity: 0.0
---
...
\end{terminal}

2 つ目のターミナル（\cmd{turtle\_teleop\_key}）で矢印キーを押してカメを動かすと、表示される値が変わることがわかります。
止めるときは \key{Ctrl}+\key{C} を押します。

\cmd{ros2 topic info} でそのトピックのパブリッシャとサブスクライバの数を、\cmd{ros2 topic hz} でデータが 1 秒間に何回届いているかを調べられます。

\begin{terminal}[トピックの情報と頻度を調べる]
$ ros2 topic info /turtle1/cmd_vel
Type: geometry_msgs/msg/Twist
Publisher count: 1
Subscription count: 1
$ ros2 topic hz /turtle1/pose
average rate: 62.502
	min: 0.015s max: 0.017s std dev: 0.00041s window: 64
...
\end{terminal}

\subsection{メッセージ型}

トピックに流れるデータには、決まった形があります。
これを\textbf{メッセージ型}といい、\cmd{パッケージ名/msg/型の名前} の形で表します。
たとえば、\cmd{/turtle1/cmd\_vel} のメッセージ型は \cmd{geometry\_msgs/msg/Twist} です。
メッセージ型の中身は、\cmd{ros2 interface show} で確かめられます。

\begin{terminal}[メッセージ型の中身を表示する]
$ ros2 interface show geometry_msgs/msg/Twist
# This expresses velocity in free space broken into its linear and angular parts.

Vector3  linear
        float64 x
        float64 y
        float64 z
Vector3  angular
        float64 x
        float64 y
        float64 z
\end{terminal}

\cmd{Twist} は、並進の速度（\cmd{linear}）と回転の速度（\cmd{angular}）を、それぞれ x・y・z の 3 つの値で表すメッセージ型です。
平面を移動するロボットでは、\cmd{linear.x}（前後の速度 [m/s]）と \cmd{angular.z}（その場での回転の速度 [rad/s]）を主に使います。
\cmd{float64} は 64 ビットの小数を表す型です。

ロボット開発では、次のようなメッセージ型がよく使われます。

\begin{tablebox}[ロボット開発でよく使うメッセージ型][tab:ros2-concepts-msgs]
  \begin{tabular}{lp{0.5\linewidth}}
    \toprule
    メッセージ型 & 内容 \\
    \midrule
    \cmd{std\_msgs/msg/String} & 文字列 \\
    \cmd{geometry\_msgs/msg/Twist} & 速度（移動ロボットへの指令など） \\
    \cmd{sensor\_msgs/msg/LaserScan} & LiDAR で測った距離のデータ \\
    \cmd{sensor\_msgs/msg/Image} & カメラの画像 \\
    \cmd{sensor\_msgs/msg/Imu} & IMU の加速度・角速度 \\
    \cmd{nav\_msgs/msg/Odometry} & 車輪の回転などから推定したロボットの位置と速度 \\
    \bottomrule
  \end{tabular}
\end{tablebox}

\subsection{コマンドからトピックにパブリッシュする}

\cmd{ros2 topic pub} を使うと、コマンドからトピックにデータをパブリッシュできます。
自分でノードを書かなくても、ほかのノードの動作を試せるので便利です。
データの中身は、YAML という形式（\cmd{\{キー: 値\}}）で書きます。

\begin{terminal}[カメに速度の指令を 1 回送る]
$ ros2 topic pub --once /turtle1/cmd_vel geometry_msgs/msg/Twist "{linear: {x: 2.0, y: 0.0, z: 0.0}, angular: {x: 0.0, y: 0.0, z: 1.8}}"
publisher: beginning loop
publishing #1: geometry_msgs.msg.Twist(linear=geometry_msgs.msg.Vector3(x=2.0, y=0.0, z=0.0), angular=geometry_msgs.msg.Vector3(x=0.0, y=0.0, z=1.8))
\end{terminal}

カメが少しだけ前に進みながら曲がります。
\cmd{--once} の代わりに \cmd{--rate 1} を付けると、1 秒に 1 回、止めるまで送り続けるので、カメは円を描いて動き続けます。

\begin{terminal}[カメに速度の指令を送り続ける]
$ ros2 topic pub --rate 1 /turtle1/cmd_vel geometry_msgs/msg/Twist "{linear: {x: 2.0}, angular: {z: 1.8}}"
\end{terminal}

省略した値は 0 になるので、このように必要な値だけを書いても構いません。
確認したら、\key{Ctrl}+\key{C} で止めましょう。

\section{サービス}
\label{sec:ros2-concepts-service}

\subsection{サービスとは}

\textbf{サービス}は、\textbf{要求（リクエスト）を送ると、応答（レスポンス）が 1 回返ってくる}通信の方法です。
サービスを提供するノードを\textbf{サーバ}、サービスを呼び出すノードを\textbf{クライアント}といいます。
お店での注文にたとえると、注文（リクエスト）をすると、品物（レスポンス）が渡されるような関係です。

トピックが「流れ続けるデータ」に向いているのに対して、サービスは「この処理をしてほしい」「この値を教えてほしい」といった、\textbf{1 回で終わるやり取り}に向いています。
ただし、クライアントは応答が返ってくるまで待つことになるので、時間のかかる処理には向いていません（その場合は、次の節のアクションを使います）。

\subsection{サービスを調べて呼び出す}

サービスの一覧は、\cmd{ros2 service list} で表示できます。
\cmd{-t} を付けると、サービスの型も表示されます。

\begin{terminal}[サービスの一覧を表示する]
$ ros2 service list -t
/clear [std_srvs/srv/Empty]
/kill [turtlesim/srv/Kill]
/reset [std_srvs/srv/Empty]
/spawn [turtlesim/srv/Spawn]
/turtle1/set_pen [turtlesim/srv/SetPen]
/turtle1/teleport_absolute [turtlesim/srv/TeleportAbsolute]
...
\end{terminal}

サービスの型の中身も、\cmd{ros2 interface show} で確かめられます。
\cmd{---} の上がリクエスト、下がレスポンスの形です。

\begin{terminal}[サービスの型の中身を表示する]
$ ros2 interface show turtlesim/srv/Spawn
float32 x
float32 y
float32 theta
string name # Optional.  A unique name will be created and returned if this is empty
---
string name
\end{terminal}

\cmd{/spawn} は、指定した位置に新しいカメを出すサービスです。
\cmd{ros2 service call} で、コマンドからサービスを呼び出してみましょう。

\begin{terminal}[新しいカメを出すサービスを呼び出す]
$ ros2 service call /spawn turtlesim/srv/Spawn "{x: 2.0, y: 2.0, theta: 0.2, name: ''}"
requester: making request: turtlesim.srv.Spawn_Request(x=2.0, y=2.0, theta=0.2, name='')

response:
turtlesim.srv.Spawn_Response(name='turtle2')
\end{terminal}

turtlesim のウィンドウに、2 匹目のカメが現れます。
レスポンスとして、新しいカメの名前（\cmd{turtle2}）が返ってきています。
ほかにも、\cmd{/clear}（カメの軌跡を消す）などのサービスを試してみましょう。

\begin{terminal}[カメの軌跡を消すサービスを呼び出す]
$ ros2 service call /clear std_srvs/srv/Empty
\end{terminal}

\section{アクション}
\label{sec:ros2-concepts-action}

\subsection{アクションとは}

\textbf{アクション}は、\textbf{時間のかかる処理}を頼むための通信の方法です。
クライアントがサーバに\textbf{目標（ゴール）}を送ると、サーバは処理を進めながら\textbf{途中経過（フィードバック）}を何度も送り返し、処理が終わったら\textbf{結果（リザルト）}を返します。
また、処理の途中で\textbf{キャンセル}することもできます。

たとえば、「目的地まで移動して」という指示を出すと、ロボットは移動しながら「残りあと 3 m」のように途中経過を知らせ、着いたら「到着した」と結果を返す、といった使い方です。
\ref{chap:practice}章で使う Navigation2 の自律移動も、アクションで指示を出します。

\subsection{アクションを調べて呼び出す}

アクションの一覧は、\cmd{ros2 action list} で表示できます。

\begin{terminal}[アクションの一覧を表示する]
$ ros2 action list -t
/turtle1/rotate_absolute [turtlesim/action/RotateAbsolute]
\end{terminal}

\cmd{/turtle1/rotate\_absolute} は、カメを指定した向きまで回転させるアクションです。
アクションの型は、\cmd{---} で区切られた 3 つの部分（ゴール・リザルト・フィードバック）でできています。

\begin{terminal}[アクションの型の中身を表示する]
$ ros2 interface show turtlesim/action/RotateAbsolute
# The desired heading in radians
float32 theta
---
# The angular displacement in radians to the starting position
float32 delta
---
# The remaining rotation in radians
float32 remaining
\end{terminal}

\cmd{ros2 action send\_goal} でゴールを送ってみましょう。
\cmd{--feedback} を付けると、途中経過も表示されます。
向き（\cmd{theta}）はラジアンで指定します（1.57 ラジアンは約 90 度です）。

\begin{terminal}[カメを回転させるアクションのゴールを送る]
$ ros2 action send_goal /turtle1/rotate_absolute turtlesim/action/RotateAbsolute "{theta: 1.57}" --feedback
Waiting for an action server to become available...
Sending goal:
     theta: 1.57

Goal accepted with ID: f8db8f44410849eaa93d3feb747dd444

Feedback:
    remaining: -1.5520000457763672
...
Result:
    delta: -1.5679999589920044

Goal finished with status: SUCCEEDED
\end{terminal}

カメがゆっくり回転し、その間、残りの回転量（\cmd{remaining}）がフィードバックとして届き、最後に結果が表示されます。
実は、\cmd{turtle\_teleop\_key} の \key{G} \key{B} \key{V} \key{C} \key{D} \key{E} \key{R} \key{T} キーも、このアクションを使ってカメの向きを変えています。

\subsection{通信の方法の使い分け}

3 つの通信の方法の違いを、表にまとめておきます。

\begin{tablebox}[トピック・サービス・アクションの使い分け][tab:ros2-concepts-compare]
  \begin{tabular}{lp{0.28\linewidth}p{0.4\linewidth}}
    \toprule
    方法 & やり取り & 向いている場面 \\
    \midrule
    トピック & 一方向に流し続ける & センサのデータ、速度の指令など、次々と流れるデータ \\
    サービス & 要求に対して 1 回応答する & 設定の変更や値の問い合わせなど、すぐに終わる処理 \\
    アクション & 目標・途中経過・結果。キャンセルもできる & 目的地への移動やアームの動作など、時間のかかる処理 \\
    \bottomrule
  \end{tabular}
\end{tablebox}

\section{パラメータ}
\label{sec:ros2-concepts-param}

\subsection{パラメータとは}

\textbf{パラメータ}は、ノードが持っている設定値です。
たとえば、ロボットの最高速度や、センサのデバイスファイルの名前（\cmd{/dev/ttyUSB0} など）を、プログラムに直接書き込むのではなくパラメータにしておけば、プログラムを書き換えずに、起動するときや実行中に設定を変えられます。
turtlesim のノードは、背景の色をパラメータとして持っています。

\subsection{パラメータを調べて変更する}

パラメータの一覧は、\cmd{ros2 param list} で表示できます。

\begin{terminal}[パラメータの一覧を表示する]
$ ros2 param list
/teleop_turtle:
  qos_overrides./parameter_events.publisher.depth
  ...
  scale_angular
  scale_linear
  use_sim_time
/turtlesim:
  background_b
  background_g
  background_r
  ...
  use_sim_time
\end{terminal}

\cmd{ros2 param get} でパラメータの値を表示し、\cmd{ros2 param set} で値を変更できます。

\begin{terminal}[パラメータの値を表示・変更する]
$ ros2 param get /turtlesim background_r
Integer value is: 69
$ ros2 param set /turtlesim background_r 150
Set parameter successful
\end{terminal}

turtlesim のウィンドウの背景の色が変わります。

\subsection{パラメータをファイルに保存する}

\cmd{ros2 param dump} を使うと、ノードのパラメータを YAML 形式で表示できます。
リダイレクト（\ref{sec:text-pipe-redirect}節）でファイルに保存しておけば、次にノードを起動するときに \cmd{--params-file} で読み込んで、同じ設定で起動できます。

\begin{terminal}[パラメータをファイルに保存して、起動時に読み込む]
$ ros2 param dump /turtlesim > turtlesim.yaml
$ ros2 run turtlesim turtlesim_node --ros-args --params-file turtlesim.yaml
\end{terminal}

パラメータをファイルで管理する方法は、\ref{chap:launch}章で詳しく学びます。

\section{rqt と rqt\_graph による可視化}
\label{sec:ros2-concepts-rqt}

\subsection{rqt\_graph}

ノードとトピックの数が増えてくると、コマンドの出力だけでは、全体のつながりがわかりにくくなります。
\textbf{rqt\_graph} を使うと、動いているノードとトピックのつながりを図で表示できます。

\begin{terminal}[rqt\_graph を起動する]
$ rqt_graph
\end{terminal}

\begin{figure}[tb]
  \centering
  \includegraphics[width=0.9\linewidth]{ros2_concepts/rqt_graph.png}
  \caption{rqt\_graph でノードとトピックのつながりを表示した例}
  \label{fig:ros2-concepts-rqt-graph}
\end{figure}

図\ref{fig:ros2-concepts-rqt-graph}のように、楕円がノード、四角がトピックを表し、矢印がデータの流れる向きを表します。
\cmd{/teleop\_turtle} から \cmd{/turtlesim} へ、\cmd{/turtle1/cmd\_vel} トピックでデータが流れていることが一目でわかります。
\cmd{/turtle1/rotate\_absolute} の中の四角は、\ref{sec:ros2-concepts-action}節で学ぶアクションが内部で使っているトピックです。
自分で作ったノードが思ったとおりにつながっているかを確かめるときに、とても便利です。
表示が更新されないときは、左上の更新ボタンを押します。

\subsection{rqt}

\textbf{rqt} は、ROS 2 のさまざまな GUI のツールを、プラグインとしてまとめて使える道具です。
\cmd{rqt} コマンドで起動し、メニューの「Plugins」から使いたいツールを選びます。
rqt\_graph も、rqt のプラグインの 1 つです。

\begin{tablebox}[rqt のよく使うプラグイン][tab:ros2-concepts-rqt]
  \begin{tabular}{lp{0.55\linewidth}}
    \toprule
    プラグイン & できること \\
    \midrule
    Node Graph & ノードとトピックのつながりを表示する（rqt\_graph と同じ） \\
    Topic Monitor & トピックに流れているデータを一覧で表示する \\
    Message Publisher & GUI でトピックにデータをパブリッシュする \\
    Service Caller & GUI でサービスを呼び出す \\
    Plot & トピックの数値をグラフにする \\
    Console & ノードが出力したログを表示する（\ref{chap:debug}章） \\
    \bottomrule
  \end{tabular}
\end{tablebox}

\begin{point}[ros2 コマンドのまとめ]
  この章で使った \cmd{ros2} コマンドは、\cmd{ros2 対象 操作} という共通の形をしています。
  \begin{itemize}
    \item 対象：\cmd{node}, \cmd{topic}, \cmd{service}, \cmd{action}, \cmd{param}, \cmd{interface}
    \item 操作：\cmd{list}（一覧）、\cmd{info}（詳細）、\cmd{echo}（表示）、\cmd{pub}（送る）、\cmd{call}（呼び出す）、\cmd{get} / \cmd{set}（取得・変更）など
  \end{itemize}
  使い方を忘れたら、\cmd{ros2 topic --help} のように \cmd{--help} を付けて調べましょう（\ref{sec:terminal-help}節）。
\end{point}

\section*{この章のまとめ}

\begin{itemize}
  \item ノードは ROS 2 のプログラムの単位です。\cmd{ros2 node list} と \cmd{ros2 node info} で調べられます。
  \item トピックは、パブリッシャからサブスクライバへデータを流し続ける通信の方法です。\cmd{ros2 topic list} / \cmd{echo} / \cmd{pub} で調べたり、データを送ったりできます。データの形はメッセージ型で決まっていて、\cmd{ros2 interface show} で確かめられます。
  \item サービスは要求に対して 1 回応答する通信、アクションは時間のかかる処理を頼み、途中経過と結果を受け取る通信です。
  \item パラメータはノードの設定値で、\cmd{ros2 param get} / \cmd{set} で表示・変更できます。
  \item rqt\_graph でノードとトピックのつながりを図で確かめられます。rqt には、ほかにも便利なプラグインがあります。
\end{itemize}
